% Fichier: listefigures.tex
% Liste des figures pour le mémoire de détection d'intrusions

\newpage
\thispagestyle{empty}

\begin{center}
    \vspace{2cm}
    {\Large \textbf{\color{blueuniv}LISTE DES FIGURES}}
    \vspace{1cm}
\end{center}

\begin{enumerate}
    \item[\textbf{Figure 1}] Architecture générale du système de détection d'intrusions hiérarchique \dotfill \pageref{fig:pipeline_general}
    
    \item[\textbf{Figure 2}] Distribution des classes dans le jeu de données NSL-KDD \dotfill \pageref{fig:data_distribution}
    
    \item[\textbf{Figure 3}] Architecture de l'autoencodeur pour la détection d'anomalies \dotfill \pageref{fig:autoencoder}
    
    \item[\textbf{Figure 4}] Illustration du processus SVM-SMOTE pour la génération d'échantillons synthétiques \dotfill \pageref{fig:smote}
    
    \item[\textbf{Figure 5}] Distribution des 10 principaux types de protocoles dans NSL-KDD \dotfill \pageref{fig:protocol_dist}
    
    \item[\textbf{Figure 6}] Comparaison des performances F1-score pour la classification binaire \dotfill \pageref{fig:binary_perf}
    
    \item[\textbf{Figure 7}] Évolution des pertes d'entraînement et de validation de l'autoencodeur \dotfill \pageref{fig:ae_training}
    
    \item[\textbf{Figure 8}] Distribution des erreurs de reconstruction pour la détection d'anomalies \dotfill \pageref{fig:reconstruction_dist}
    
    \item[\textbf{Figure 9}] Impact de SVM-SMOTE sur la distribution des classes \dotfill \pageref{fig:smote_impact}
    
    \item[\textbf{Figure 10}] Taux de réussite par classe dans la classification multi-classes \dotfill \pageref{fig:class_perf}
    
    \item[\textbf{Figure 11}] Matrice de confusion normalisée pour la classification multi-classes \dotfill \pageref{fig:confusion_matrix}
    
    \item[\textbf{Figure 12}] Métriques détaillées par classe : Précision, Rappel et Spécificité \dotfill \pageref{fig:detailed_metrics}
    
\end{enumerate}

\vfill

\begin{center}
    \textit{Total : 12 figures}
\end{center}

\newpage